% Section 6 (owner: phil). Plan: ../OUTLINE.md, section "sec:phil"; rows: ../CLAIMS.md §5; notation: ../NOTATION.md.
% Sources: philosophy §1-§3 (final record), synthesis §2.1-§2.5 and §2.8 (as corrected by final-check §4).
% Every philosophical sentence is an assessment argued in philosophy §1-§3; sentences that rest on a theorem cite it.
% For the page budget, prop:phil:conservation, prop:phil:predicting and prop:phil:flagpolesystem are stated in App. D (app-phil.tex) and
% summarised here in prose (logged in ../logs/write-phil.md).
% Local symbols (this section and App. D only): [k]; Un_D(Q); the conservation world m, Z, sigma, t^*, kappa_t, t_max and
% the predictors Raw^+, Raw^-, Qa^+, Qa^-, Qa^bin, Mix; the raw alphabet \mathbb X, mu, eta, mu_eta, f_x, (phi_q);
% the flagpole ht, sh, u, p_i, a_i, N_p, Lambda_h, Lambda_s, T_H, T_S, T_H', T_S', C_H, C_S;
% in App. D also Delta_sh, sh_H, ht_S, qt_j, theta_j, Lambda_j^ht, Lambda_j^sh, Delta_split, z, Sunlit, p'_l, a'_l, N_d, Psi.
\section{What a theory does and what it is}\label{sec:phil}

What does a scientific theory do? It may predict observations, the view behind Solomonoff induction \citep{solomonoff1964formal}; answer questions posed in some language, why-questions included \citep[ch.~5]{vanfraassen1980scientific}; explain \citep{hempel1948studies}; or solve problems \citep{laudan1977progress}. And what is it? A program that predicts observations; a set of sentences closed under consequence \citep{suppe1977structure}; a class of models \citep{suppes1960comparison,vanfraassen1980scientific}; a core of models with intended applications \citep{sneed1971logical,balzer1987architectonic}; or a practice of exemplars and skills \citep{kuhn1970structure,kuhn1977second,polanyi1958personal}. The \emph{axiom view} answers that a theory answers questions posed in a given language and is a logical theory; the \emph{raw-data view}, that it predicts observations and is a program. We assess the axiom view against the theorems of this paper, and say which of our claims are theorems (proofs in \cref{app:phil}) and which are arguments.

\subsection{An argument map}\label{sec:phil:map}

We separate five claims (\cref{tab:phil:map}). \textbf{D}: a theory answers questions posed in a given language rather than predicting raw data, read comparatively, as a step away from raw-data prediction rather than the end point. \textbf{S}: a theory is a logical theory, not a program. \textbf{E}: theories explain propositions. \textbf{L}: the ladder of \cref{sec:intro} (raw data; questions posed with existing abstractions; theories that extend the language; theories that revise the base abstractions). \textbf{P}: \SAIabbr is progress over raw-data prediction, both in what a theory does and in what it is.

\begin{table}[t]
\crefabbrev
\centering\footnotesize
\renewcommand{\arraystretch}{1.08}
\begin{tabularx}{\textwidth}{@{}l X p{0.31\textwidth}@{}}
\toprule
claim & our assessment (an argument) & formal results it rests on \\
\midrule
\textbf{D} answering & true in comparative form; not a new inference principle; survives the objection from statistical practice & \cref{thm:soi:coinfree,thm:soi:onetape,prop:phil:answers,prop:phil:conservation,prop:phil:extractor,prop:phil:predicting,prop:soi:closure,prop:soi:independent} \\
\textbf{S} logical form & descriptive part true of declarative textbook presentations (an observation); explanatory part open, and relative to a reasoner & \cref{sec:phil:craig}; \cref{thm:shape:invisible,thm:shape:directions,cor:shape:regimes,thm:terms:dilemma} \\
\textbf{E} explanation & true in the unification sense, with at least four gaps & \cref{prop:phil:flagpole,prop:phil:flagpolesystem,prop:laws:correction,prop:laws:copernicus} \\
\textbf{L} the ladder & \SAIabbr is on step 2, with partial results on step 3; step 4 has no formal handle & \cref{thm:terms:horn,thm:terms:dilemma,cor:shape:regimes} \\
\textbf{P} progress & to be replaced by a list of formal differences; none concerns the shape of hypotheses in the idealised setting & \cref{tab:phil:differences} \\
\bottomrule
\end{tabularx}
\caption{The argument map: our assessments, which are arguments, and the results they rest on, which are theorems in their stated settings.}
\label{tab:phil:map}
\end{table}

\paragraph{What rests on theorems.}\status{assessment}\src{philosophy §1.1-§1.7; syn §2.2, §2.8}
Proved, in the stated settings: the idealised formalism cannot see the difference between a theory and a program (\cref{sec:phil:craig}); answering questions differs from raw-data prediction in target, data type, coherence and partiality, not in the inference principle (\cref{sec:phil:questions}); with a cost on the written theory, usefulness depends exponentially on the pair (reasoner, vocabulary), and on declarative against procedural form only by a constant factor once intermediate concepts are available and the reasoner can chain through them (in the model of \cref{sec:shape}; open beyond it); theoretical vocabulary is dispensable in the ideal mixture, while relations or entities pay for finite theories and under charged deduction (\cref{sec:terms}); and ``Kepler's laws are only approximately right'' is a limitation of proof-requiring compatibility with deterministic theories, not of the axiom view, since in $\MA$ an approximately true law is a law with a calibrated error model (\cref{sec:laws}). Computed, conditional on the pool of forms and the codes: on noise-free interval data, simplicity among compatible forms is relative to the language, and with graded fit the data decide (\cref{sec:kepler}). A theorem plus an argued identification: Hempel's inductive-systematisation reply fails in the idealised setting (\cref{rem:phil:hempel}). Arguments only: the descriptive observation; why declarative shape is good (open); that D survives the objection from statistical practice; the readings of the toys and the general claims about explanation; that the axiom view misses the procedural, statistical and dynamic sides of theories (\cref{sec:limits}); and our refusal, resting on a possibility result (\cref{cor:logic:false}), to say anything realist about induced theories.

\paragraph{The ladder.}\status{assessment}\src{philosophy §1.6}
We accept L. The base language of step 2 may be highly theoretical (heliocentric longitudes presuppose Copernican geometry); no theory-neutral observation language is assumed, so the classical objections to one do not apply in the same way. But theory-ladenness \citep{hanson1958patterns} biases comparisons between theories that disagree about the base language, since a theory that revises it is scored in a language it rejects; this is why the ladder has a step 4.

\subsection{The descriptive observation and the explanatory question}\label{sec:phil:observation}

Claim S combines a description with an explanation, which we keep apart.\status{assessment (observation); open}\src{philosophy §1.2, §1.4; syn §2.3}
\emph{The observation:} textbook presentations of fundamental theories are declarative. They state laws, definitions, principles and models given by equations, not procedures that print predictions; this needs no theory of why it holds. Two qualifications. \emph{Declarative is not sentential:} physics texts write something closer to a class of models specified by equations \citep{suppes1960comparison,suppes1962models} than to first-order axioms. We use sentences because questions and data are sentences and the prior of \SAIabbr is over texts; the question-answering view itself is neutral between sentences and classes of models, since in partition semantics a question is a partition of logical space (cf.\ \citealp{groenendijk1984studies}). \emph{Theories in use have procedural parts:} methods attached to concepts, numerical schemes, simulation codes; in quantum chemistry a ``level of theory'' is a computational method, and one position of Mars from Kepler's laws already needs root-finding for Kepler's equation. So the observation concerns textbook presentations, not theories in all their uses.

\emph{The question}, why the declarative shape is good and whether an ideal inducer should produce it, is open. Craig's construction rules out an answer based only on what can be represented (\cref{sec:phil:craig}), so an answer must be relative to a reasoner with costs, or to a community of reasoners. In the model of \cref{sec:shape}, the best formal answer we have, what matters is a reasoner that can do algebra or chain through intermediate concepts, together with those concepts, not logical form as such. The explanation may lie in communication and piecemeal checking, which have no formal setting yet.

\begin{question}[Why declarative shape?]\label{q:phil:declarative}
\status{open}\src{syn §6 item 26; philosophy §1.2, §5 (L11)}
Is there a formal setting for a community of bounded reasoners who communicate, check and extend theories piece by piece, in which theories of declarative shape provably do better than programs?
\end{question}

\status{assessment, resting on thm:shape:invisible, thm:shape:directions (proved)}\src{philosophy §1.4 (correction of S4); syn §2.3}A natural argument for S is that a theory answers in many directions (Kepler's third law gives the period from the semi-major axis and conversely; \cref{def:kepler:language}), while a program computes one direction cheaply. The second half holds of particular programs, not of programs as a class: a generic solver applied to a declarative law costs $O(1)$ extra length, a program with internal variables answers all $m$ directions of a parity law on $m$ variables with $O(m)$ operations, and under program length directions cost nothing. They matter only for a fixed weak reasoner and a cost on the written theory: for one parity law, forward chaining needs exactly $m$ times as many rules as unit propagation needs clauses, over the base variables or with one round, and at most $9$ times the written length of the least clausal theory with intermediate concepts and unbounded rounds (\cref{thm:shape:invisible,thm:shape:directions}).

\subsection{What a theory does: questions against raw data}\label{sec:phil:questions}

Questions, and the answers a theory proves, are defined in \cref{def:setting:questions}, following the logic of questions (cf.\ \citealp{belnap1976logic,hintikka1999inquiry}); $[k]:=\{1,\dots,k\}$.

\begin{definition}[Answers]\label{def:phil:questions}
\status{definition}\src{philosophy Def 2.0}
A theory $T$ \emph{answers} a question $\mathcal Q=\{\varphi_1,\dots,\varphi_k\}$ if $B\cup T\vdash\varphi_i$ for some $i$, and \emph{answers it partially} if $B\cup T\vdash\bigvee_{i\in J}\varphi_i$ for some $J$ with $\emptyset\ne J\subsetneq[k]$.
\end{definition}

In the Kepler example (\cref{def:kepler:language}) a question is a cover by interval sentences (``where is Mars at time $t$, to within $2'$?'') or structural (``is the orbit of Mars planar?''); without the elements of the orbits, Kepler's laws answer structural questions but entail no direction datum with tolerance below $\pi$ (\cref{prop:kepler:ramsey}). \status{assessment}\src{philosophy §2.1 reading}A raw-data predictor answers one question, ``what is the next observation?'', and answers it fully. \SAIabbr answers every question the base language can pose, and reports how much of its posterior answers none:

\begin{proposition}[The answer distribution]\label{prop:phil:answers}
\status{proved; computed}\src{philosophy Prop 2.1}
Let $\MA(D)>0$, let $\mathcal Q=\{\varphi_1,\dots,\varphi_k\}$ be a question, and let $\mathrm{Un}_D(\mathcal Q):=\Unif(E_D\setminus\bigcup_iO_{\varphi_i})/\Unif(E_D)$, the posterior probability that the random theory proves no direct answer. (a) $\sum_{i\le k}\Bel_D(\varphi_i)+\mathrm{Un}_D(\mathcal Q)=1$. (b) For $\emptyset\ne J\subsetneq[k]$, $\Bel_D(\bigvee_{i\in J}\varphi_i)\ge\sum_{i\in J}\Bel_D(\varphi_i)$, and the difference is the posterior probability that the random theory answers $\mathcal Q$ partially through $J$ but proves no $\varphi_i$ with $i\in J$; for $J=[k]$, $\Bel_D(\bigvee_i\varphi_i)=1$ and the difference is $\mathrm{Un}_D(\mathcal Q)$. (c) For $\mathcal Q=\{\varphi,\neg\varphi\}$, (a) is the trichotomy $\Bel_D(\varphi)+\Dis_D(\varphi)+\Indep_D(\varphi)=1$ of \cref{prop:basic:structure}.
\end{proposition}

On $E_D$ the random theory is consistent, so it proves at most one direct answer (\cref{app:phil:answers}). In \SOIabbr the analogue of $\mathrm{Un}_D$ mixes logical independence with a labeller's silence (\cref{prop:soi:closure}).

\paragraph{The conservation-law world.}\status{assessment, resting on prop:phil:conservation (proved; computed)}\src{philosophy Prop 2.2, §2.2 reading}
What is a law worth to each task? In a toy world (\cref{prop:phil:conservation} in \cref{app:phil:conservation}) a micro-state $x\in\{0,1\}^m$ moves by a uniformly random permutation of each weight class, so its weight is conserved; the raw data are the trajectory, and the questions are ``what is $\lvert x_t\rvert$?''. Knowing the law saves a Bayes predictor of the raw data between $\frac12\log\frac{\pi m}2+o(1)$ bits and one bit more per step (for $t\le Z/2$, $Z$ the size of the weight class), a share $O(\log m/m)$ of its step loss, until the trajectory returns to its start; on $t_{\max}+1$ questions it saves the share $t_{\max}/(t_{\max}+1)$ of the no-law loss. Yet by the chain rule the raw saving at each step equals the loss of the raw no-law predictor's own answer to ``what is $\lvert x_{t+1}\rvert$?'', and in each task a Bayes mixture of the two hypotheses pays at most one bit for learning the law. A computation (code in the accompanying repository) gives savings of $2.348$, $3.331$ and $5.326$ bits at the start of a trajectory at $m=16$, $64$, $1024$, against $\frac12\log\frac{\pi m}2=2.326$, $3.326$, $5.326$, and shares $0.15$, $0.052$, $0.0052$. So the share of the raw loss that the law saves tends to $0$, because of micro-detail that no theory addresses, while its absolute gain per step equals its value for the corresponding question, and both tasks learn the law at the same price: the theory is ``in'' the raw predictor. The tasks differ in their \emph{target}, which losses are counted, not in what must be learned nor in what the law is worth per step; scoring questions is right when what we want is answers. (The saving on questions holds by construction: the toy asks only about the conserved quantity.)

Answers never cost more description than raw modelling plus an \emph{extractor} that reads them off the raw data. Let $\mathbb X$ be a finite alphabet, $\mu$ a computable probability measure on $\mathbb X^{\mathbb N}$ (the raw world), $(\varphi_q)_{q\in\mathbb N}$ a fixed computable family of sentences (the questions), and $\eta$ a partial computable function with $\eta(q,x)\in\{0,1\}$ computed from a finite prefix of $x$; put $f_x(\varphi_q):=\eta(q,x)$ and $\mu_\eta(D):=\mu\{x:f_x\supseteq D\}$.

\begin{proposition}[Raw modelling plus an extractor]\label{prop:phil:extractor}
\status{proved, given thm:basic:dominance and the standard sampling construction}\src{philosophy Prop 2.4}
If $f_x$ is consistent with $B$ for $\mu$-almost every $x$, then $\MA(D)\ge2^{-K(\mu)-K(\eta)-c}\mu_\eta(D)$ for every finite data set $D$ of question--answer pairs, so $\LA(D)\le-\log\mu_\eta(D)+K(\mu)+K(\eta)+c$; here $c$ depends also on the family $(\varphi_q)$. For a computable sequence $x$ with $f_x$ consistent with $B$, $\LA(D)\le K(x)+K(\eta)+c$ on every $D$ with $f_x\supseteq D$.
\end{proposition}

A theory machine samples the raw world and emits the extractor's labels (\cref{app:phil:extractor}). The extractor is the bridge from raw data to the question language, and its cost $K(\eta)$ is not small in real cases: it contains whatever recognises a planet in an image. Answering is also the stronger achievement (\cref{prop:phil:predicting} in \cref{app:phil:predicting}): a consistent r.e.\ theory that proves the data bounds the loss of $\MA$ by its complexity, but on the lottery stream $\MA$ predicts well while no consistent r.e.\ theory proves the data; predicting data that present an incomplete theory leaves its open questions open; and the beliefs of \SAIabbr are monotone under entailment and closed under conjunction up to additive error, unlike, in the semimeasure convention, the predictions of consistent \SOIabbr (\cref{prop:soi:closure}).

\paragraph{The objection from statistical practice.}\label{rem:phil:practice}\status{assessment}\src{philosophy §1.3 (o4), §2.6; syn §2.2 row D}
Much of physics confronts theory with near-raw data through generative models and information-theoretic losses: gravitational-wave searches match waveform templates against the strain series with a Gaussian noise likelihood, cosmology fits likelihoods of microwave-background data, particle physics simulates detectors event by event. D survives in comparative form, we think, because these pipelines exist to answer a few questions (a mass, a spin, a cosmological parameter, whether a particle exists). But the generative model is theory plus instrument plus noise, and the proof-requiring core of \SAIabbr, with no native noise model, sits far from this practice (\cref{sec:limits:stat}). So D holds in comparative form; its formal content is a change of data type, logical coherence, partiality and the choice of target, not a new inference principle (one-tape \SAIabbr $\asymp$ one-tape consistent \SOIabbr, \cref{thm:soi:onetape}); and the real difficulty, the extractor, \SAIabbr takes as given.

\paragraph{Purpose and logarithmic loss.}\label{rem:phil:purpose}\status{assessment}\src{philosophy §2.2 remark; syn §7 item 2}
Grading a prediction for use and scoring credences are different jobs. For the first, the purpose decides what is graded and with which tolerance, as with the interval data of \cref{sec:kepler}; \SAI lets the purpose fix the propositions, and scores credences on those. For the second, the logarithmic score is, up to positive affine transformation, the only strictly proper local scoring rule on more than two outcomes (a result attributed to \citealp{bernardo1979expected}), but which proper score is privileged is contested (accuracy-first arguments favour quadratic inaccuracy, \citealp{leitgeb2010objective,pettigrew2016accuracy}), and for bounded or misspecified models different proper losses select different best models. We do not claim that the logarithmic score is privileged; even inside \SAIabbr the price of a datum can depend on the loss convention by an unbounded amount (\cref{prop:logic:normalised}).

\subsection{What a theory is, after Craig}\label{sec:phil:craig}

\paragraph{The idealised formalism cannot see shape.}\label{rem:phil:craig}\status{assessment, resting on thm:soi:coinfree, thm:soi:onetape, cor:basic:vocabulary, thm:terms:horn, thm:terms:KCV, thm:shape:invisible, prop:phil:flagpolesystem, prop:soi:freededuction (proved)}\src{syn §2.8 row 1; philosophy §1.4 item 1}
Every consistent assigner of truth values is an axiom set at $O(1)$ extra bits. Coin-free \SAIabbr $\asymp$ coin-free consistent \SOIabbr, by Craig's construction inside the universal prior (\cref{thm:soi:coinfree}), and one-tape \SAIabbr $\asymp$ one-tape consistent \SOIabbr (\cref{thm:soi:onetape}). Adding vocabulary changes the weights of old-vocabulary data by $O(1)$ (\cref{cor:basic:vocabulary,thm:terms:horn}); coin-free, r.e.\ axiom induction is finite-theory induction in a three-sorted extension, up to $O(1)$ (\cref{thm:terms:KCV}). Program length sees neither the reasoner nor the vocabulary (\cref{thm:shape:invisible}); the posterior cannot move between deductively equivalent formulations, whose complexities in the flagpole toys below differ by $O(1)$ however many laws are added (\cref{prop:phil:flagpolesystem}); and the diagonal separations from \SOIabbr under time bounds come from free deduction, not from axioms (\cref{prop:soi:freededuction}). So the idealised formalism cannot see the difference between a theory and a program. Nor does D support S as a claim about what an ideal inducer should produce: the labellers of \SOIabbr carry no theory and are equivalent to theories, and \citet{khardon1997learning} give a learning-to-reason algorithm for a class of propositional languages that has no efficient reasoning algorithm when represented as a formula-based knowledge base.

Where shape matters, it is relative to a reasoner and a vocabulary, and only when the written theory is charged: in the model of \cref{sec:shape} the exponential separations are between pairs (reasoner, vocabulary) (\cref{cor:shape:regimes}); in the settings of \cref{sec:terms} what pays is theoretical relations and entities, not vocabulary (\cref{thm:terms:dilemma}). We still find it plausible that good theories are shaped like logical theories; what is missing is a formal setting in which the shape matters (\cref{q:phil:declarative}), and nothing in the formal results favours sentences over classes of models. Since the limit posterior can live on random theories, and in true arithmetic with true but adversarially chosen data on false ones (\cref{thm:logic:lottery,cor:logic:false}; a possibility result), we say nothing realist about induced theories.

\subsection{Explanation}\label{sec:phil:explanation}

In the axiom view $T$ explains $\varphi$ when $B\cup T\vdash\varphi$, the deductive-nomological pattern \citep{hempel1948studies}, and an explanation is better the more a short theory derives, as in unification accounts \citep{friedman1974explanation,kitcher1981explanatory}; the explanatory power of \cref{def:laws:EP} makes this precise (\cref{prop:laws:EP} relates it to Myrvold's account, \citealp{myrvold2003bayesian}). But derivability is symmetric where explanation is not: a pole's height and the sun's elevation entail the shadow's length, and the shadow's length and the elevation entail the height (Bromberger's example, \citealp{bromberger1966why}, as attributed by \citealp{glymour1978two}). Over the theory $B$ of real closed fields, with poles $p_i$, heights $\mathrm{ht}(p)$ and shadow lengths $\mathrm{sh}(p,u)$ when the tangent of the sun's elevation is $u>0$, the shadow law has two formulations: $\Lambda_h$ says $\mathrm{sh}(p,u)\cdot u=\mathrm{ht}(p)$; $\Lambda_s$ says $\mathrm{sh}(p,u)\cdot u=\mathrm{sh}(p,1)$ and $\mathrm{ht}(p)=\mathrm{sh}(p,1)$. With initial conditions on the heights or on the reference shadows $\mathrm{sh}(p_i,1)$ they give $T_H$, $T_S$ (with $\Lambda_h$) and $T'_H$, $T'_S$ (with $\Lambda_s$). A derivation is charged by the number of closed instances of the law's universally quantified conjuncts that it uses; $C_H$, $C_S$ extend $T_H$, $T_S$ by intervention vocabulary with the height, resp.\ the reference shadow, exogenous (precise forms in \cref{app:phil:flagpole}).

\begin{proposition}[The flagpole]\label{prop:phil:flagpole}
\status{proved; computed}\src{philosophy Prop 3.1}
(a) $T_H$, $T_S$, $T'_H$, $T'_S$ are deductively equivalent over $B$; so every data set is compatible with all four or with none, in the coin-free version the posterior odds of any two programs for them equal their prior odds on every data set compatible with them, and $\lvert K(T_S)-K(T_H)\rvert\le c$. (b) Under derivation charges, every shadow datum $\mathrm{sh}(p_i,u)=v$ with $u\ne1$ that they entail needs exactly one instance from $T_H$ and $T'_S$, whose initial conditions sit at the hub of their formulation, and exactly two from $T_S$ and $T'_H$: which quantity is the hub is a property of the formulation, not of the world. (c) $C_H$ and $C_S$ have the same consequences in the old vocabulary, and one interventional datum separates them.
\end{proposition}

\paragraph{What the toys show.}\status{assessment}\src{philosophy §3.2 reading, §3.3, §3.5}
The explanation relation of \SAIabbr is direction-free, and its posterior cannot move between deductively equivalent theories. Kitcher's unificationist answer \citep{kitcher1989explanatory}, in its usual summary, is that the best systematisation of all our beliefs derives the dimensions of objects from their origin and development; a height-from-shadow pattern would add a pattern, which fails for objects that cast no shadow. With $k$ further laws in which the height figures (\cref{prop:phil:flagpolesystem} in \cref{app:phil:system}), the formulation with the shadow as hub is longer by an amount growing with $k$ in a literal code without definitions (relative to the primitives, or because some objects cast no shadow), by a constant with one definition, and by $O(1)$ under the universal prior; under derivation charges only a mixed theory loses. So Kitcher's answer is an effect of whole-theory simplicity that \SAIabbr reproduces only under a listing code, like the textbook point of \cref{sec:shape}. Direction enters cleanly with interventional questions, once the base language can pose them; but the intervention semantics was written into $C_H$ and $C_S$ by hand, so the toy is consistent with the interventionist account \citep{woodward2003making} and gives it no support over rivals such as the causal-process account \citep{salmon1984scientific}.

We conclude that derivation from a short theory is explanation in the unification sense, with at least four known gaps; we do not claim that derivation is explanation outright. The gaps are \emph{asymmetry}, as above; \emph{explanation by approximate correction}, since Newton's theory does not entail Kepler's laws but, in an exact toy, proves the honest old law on every window and refutes the exact one (\cref{prop:laws:correction}), so that ``$T$ explains $\varphi$'' must include ``$T$ derives an honest approximation of $\varphi$ with an error bound on a domain''; \emph{irrelevance}, as in the standard examples of hexed salt and of a man taking birth-control pills, where the length prior penalises an irrelevant axiom in a whole theory but a derivation does not mark which axioms it used; and the \emph{conjunction problem} ($L$ is derived from $L\wedge L'$), which the length prior and the explanatory power address in part. Unification can register an explanatory gain without a gain in accuracy: against Ptolemy's system with the coincidences of retrogression and opposition stated as axioms, which has the same fit, Copernicus's wins at every steepness exactly when the heliocentric reorganisation costs fewer bits than the coincidences (\cref{prop:laws:copernicus}).

\subsection{What separates \SAIabbr from raw-data prediction}\label{sec:phil:differences}

We replace claim P by the formal differences between \SAIabbr and a predictor of an observation stream, sorted into advantages for claims D and E, neutral differences and liabilities (\cref{tab:phil:differences}).\status{assessment (the sorting); items as cited}\src{philosophy §1.7; syn §2.5}
What does not separate the two is the inference principle, a Bayes mixture over a universal class, and, in the idealised setting, the shape of hypotheses. This is the progress claim we believe: for what a theory does, \SAIabbr is a step in the right direction; for what a theory is, the idealised formalism is neutral.

\begin{table}[t]
\crefabbrev
\centering\footnotesize
\renewcommand{\arraystretch}{1.04}
\begin{tabularx}{\textwidth}{@{}r X p{0.19\textwidth}@{}}
\toprule
& difference & support \\
\midrule
\multicolumn{3}{@{}l}{\emph{Advantages for claims D and E}}\\
1 & \emph{Data type:} labelled sentences of a chosen language; questions are sets of sentences. & \cref{def:setting:questions,def:phil:questions} \\
2 & \emph{Logical coherence:} consequences of the data get posterior probability $1$, while in the semimeasure convention consistent \SOIabbr is never certain of an unasked consequence; $\Bel$ is monotone under entailment and closed under conjunction up to additive error. & \cref{prop:soi:closure,prop:phil:predicting} \\
3 & \emph{Partiality:} ``independent'' is a native third value, whose mass is the semimeasure defect; for a $k$-way question the answer masses and the unanswered mass sum to $1$. & \cref{prop:basic:structure,prop:phil:answers} \\
4 & \emph{An observed ``independent'' label} separates theories from labellers by an unbounded amount on one datum, coin-free and one-tape, when $\mathcal L\supseteq\Lar$ and $\Q\subseteq\Cn(B)$. & \cref{prop:soi:independent} \\
5 & \emph{Laws are sentences with belief values:} $\Bel_D(\forall x\,\varphi)$ differs from the probability that all future data conform. & \cref{prop:logic:omega} \\
6 & \emph{Answers cost no more description than raw modelling plus an extractor,} and can cost far less. & \cref{prop:phil:extractor,prop:phil:conservation} \\
\midrule
\multicolumn{3}{@{}l}{\emph{Neutral differences}}\\
7 & \emph{Piecemeal confirmation:} credence in single laws, though only whole theories meet the data: the standard Bayesian treatment of Duhem's problem \citep{dorling1979bayesian}, which \citet{strevens2001bayesian} argues begs the question. \SAIabbr inherits the debate; its allocation of blame is fixed by the prior, hence by the syntax. & definitions; an argument \\
8 & \emph{Truth from unprovability:} if $\Q\subseteq\Cn(B)$, belief or independent mass on a $\Pi_1$ sentence certifies its truth; a curiosity of arithmetic data. & \cref{prop:logic:shape} \\
9 & \emph{Explanatory direction needs questions the language can pose:} interventional data separate causal theories that the old vocabulary cannot. & \cref{prop:phil:flagpole}(c); in general an argument \\
\midrule
\multicolumn{3}{@{}l}{\emph{Liabilities}}\\
10 & \emph{Random theories, false theories:} on the lottery stream bounded one-tape loss is achieved by random theories and by no consistent r.e.\ theory; in true arithmetic with true (adversarial) data the limit posterior can live on false theories, a possibility result. & \cref{thm:logic:lottery,cor:logic:false} \\
11 & \emph{Language-relative simplicity, within bounds:} faithful translations change weights by $2^{\pm(K(\tau)+c)}$; within that bound the base language decides real choices, such as the preferred orbit of Mars on noise-free interval data. & \cref{thm:basic:translation}; \cref{sec:kepler} \\
\bottomrule
\end{tabularx}
\caption{What formally separates \SAIabbr from raw-data prediction. Items 7 and 9 and the sorting are arguments; the other items are results.}
\label{tab:phil:differences}
\end{table}
